% \noindent using the three envelopes ....
% - Results: coverage by class, set size distribution, singleton accuracy
% - analysis of hard cases: ambiguous gram-positive phages

% \subsubsection{Experimental Setup}
% \noindent - Data: split-0 phages only (80\% calibration, 20\% test), unknown host types excluded\\
% - Parameters: $\alpha = 0.1 , M=200$ directions (radial),
% M=20 bins (strip), 10 independent runs to estimate variance\\
% - The NC transformation used: keep scores as-is for one class, flip with 1 - s for the other 
% % (motivated by the empirical score semantics described in Chapter 4)

% \subsubsection{Results}
% \noindent - table: overall coverage, per-class coverage, average set size, empty set rate, size-2 rate, singleton accuracy std across 10 runs for all three methods\\
% - Bar charts: set size distribution, overall and per-class coverage


% \subsubsection{Analysis of Hard Cases}
% The minority class consistently has lower coverage than the majority class before the class-conditional fix — show this 
% explicitly with a before/after comparison
% Identify the ambiguous phages: those whose mean raw score falls in the range [0.3,0.7][0.3, 0.7]
% [0.3,0.7] — these are phages for which the underlying classifiers are genuinely uncertain
% Show that coverage on typical phages of each class is much higher than coverage on the ambiguous subset — this is a property 
% of the data, not a flaw in the conformal framework
% Discuss variance across runs: radial method has higher variance than strip and collapsed due to the random direction sampling


\section{Gram Type Classification}
\label{sec5.2: gram results}
We evaluated the three envelope geometries (Radial, Strip, and Collapsed 1D) on the real phage data for a binary 
Gram type classification ($\mathcal{H} = \{\text{gram-neg}, \text{gram-pos}\}$). We only used the 
$4422$ phages in split 0 to ensure that all scores remain out of fold by construction and there is no contamination.


\subsection{Experimental Setup}
\label{sec5.2.1: gram setup}
We partitioned the split 0 subset into shape discovery ($|S_1| = 1\,768$), size scaling ($|S_2| = 1\,326$), and 
test evaluation ($|S_{\text{test}}| = 1\,328$) sets using a $40/30/30$ split. For each geometry, we fitted a single shared envelope $E$ 
on $S_1$, and then calculated the class conditional scaling factors $\hat{t}_{\text{neg}}$ and $\hat{t}_{\text{pos}}$ on $S_2$ 
using \eqref{eq: class specific quantile}. To ensure validity across both classes simultaneously, we scaled the envelope by 
their maximum:
\begin{equation}
\label{eq: gram-max-scaling}
\hat{t} = \max\big(\hat{t}_{\text{neg}},\, \hat{t}_{\text{pos}}\big).
\end{equation}

At test time, we evaluated the Gram type nonconformity transformation under both candidate hypotheses. For an empty 
candidate set ($\emptyset$), we returned both labels ($C(\mathbf{s}) = \mathcal{H}$) as established in Section~\ref{sec3.5: t 
hat calculation}.

%%%%%%%%%%%%%%%%%%%%%%%%%%%%%%%%%%%%%%%%%%%%%%%%%%%%%%%%%%%%%%%%%%%%%%%%%%

\subsection{Coverage and Prediction Efficiency}
\label{sec5.2.2: gram coverage}

\begin{figure}[h]
    \centering
    \includegraphics[width=1\textwidth]{Plots/comparison_plot.png}
    \caption{Gram type classification performance across Radial, Strip, and Collapsed 1D geometries on the test split 
            (Fold 0, nominal target $1-\alpha = 0.90$). Panels show (left to right): prediction set size distribution, marginal empirical 
             coverage, class-conditional coverage, and singleton accuracy with sample counts ($n$).}
    \label{fig: gram results}
\end{figure}

As we can see in Figure~\ref{fig: gram results} (middle panels), all three envelope geometries achieve $99.8\%$ empirical 
coverage, which is well above the nominal $90\%$ target (red dashed line). The class conditional breakdown 
confirms that the coverage is balanced for both Gram negative ($99.9\%$) and Gram positive ($99.8\%$) phages. This demonstrates 
that the max threshold calibration $\hat{t} = \max(\hat{t}_{\text{neg}}, \hat{t}_{\text{pos}})$ effectively eliminates coverage 
disparity between classes.

We can also see that the high coverage does not compromise efficiency. The set size distribution chart of Figure~\ref{fig: 
gram results} (left panel), shows that $94\text{--}95\%$ of test phages receive singleton predictions, while only $5\text
{--}6\%$ require size 2 uncertain sets ($C(\mathbf{s}) = \mathcal{H}$). The empty set rate is $0\%$ by construction due to the 
conservative fallback rule. When singleton sets are produced, the prediction accuracy is $99.8\%$ across all three geometries 
($n \ge 1\,249$), which tells us that when the model is confident enough to output exactly one label, it is almost always correct. 

Acheiving this coverage that is over the target highlights a trade off, it suggests the calibrated 
envelopes are wider than strictly necessary to meet $1-\alpha=0.90$ (due to the max threshold calibration). So we are trading 
efficiency for a safety margin the task may not require. We will return to this trade off in Section~\ref{sec6: conclusion}.


%%%%%%%%%%%%%%%%%%%%%%%%%%%%%%%%%%%

\subsection{Score Separation and Geometric Equivalence}
\label{sec5.2.3: score separation}

\begin{figure}[htbp]
\centering
\includegraphics[width=0.88\linewidth]{Plots/raw_score_distribution.png}
\caption{Distribution of raw scores per test phage, separated by true Gram type. Gram negative phages 
concentrate tightly near 0 while Gram positive phages concentrate near 1, leaving the intermediate ambiguous zone 
($0.25\text{--}0.75$, dashed red lines) almost empty.}
\label{fig: gram score separation}
\end{figure}

All three envelope constructions yield nearly identical performance metrics in Figure~\ref{fig: gram results}. 
This equivalence can directly be explained by the empirical raw score distribution in Figure~\ref{fig: gram score separation}: 
Gram negative phages concentrate tightly near the origin (mean $0.055 \pm 0.054$), whereas Gram positive phages concentrate 
near $1.0$ (mean $0.918 \pm 0.066$). The two classes form disjoint, well separated clusters with negligible density 
in the ambiguous zone ($0.25\text{--}0.75$), any sensible star shaped boundary cleanly separates them. 

This near identical 
performance across the three geometries suggests that the choice of envelope geometry is less critical for low dimensional task 
with strongly separated scores. The advantages of non convex can be seen more prominently in the 
higher dimensional, multi class host level classification task in Section~\ref{sec5.3: host results}
%%%%%%%%%%%%%%%%%%%%%%%%%%%%%%%%%%%%%%%%%%%%%%%%%%%%%%%%%%%%%%%%%%%%%%%%%%%%%%%%%%%%%%%%%%%%%%%%%%%%%%%%%%%%%%%%%%%%%




